\begin{table*}[t]
\centering
{\small
\resizebox{0.98\textwidth}{!}{
\begin{tabular}{@{} >{\raggedright\arraybackslash}p{3cm} >{\raggedright\arraybackslash}p{12cm} @{}}
\toprule
\textbf{Original sentence} & The \textit{added constraint is important}, as it \textit{prevents the model from} deviating too far from the distribution on which the reward model is accurate, as well as \textit{maintaining the generation diversity} and preventing mode-collapse to single high-reward answers. \\
\midrule
\textbf{Factual probe} & In the paper `Direct Preference Optimization: Your Language Model is Secretly a Reward Model', the authors state that the \textit{added constraint} in the reinforcement learning objective not only \textit{prevents the model from} deviating too far from the distribution on which the reward model is accurate, but also \textit{maintains} \textbf{the generation diversity}. \\
\midrule
\textbf{Related sentences} & Even if we use the MLE estimate $r_{\phi}$ of the ground-truth reward function $r^*$, it is still \textit{expensive to estimate the partition function} $Z(x)$, which makes this representation hard to utilize in practice (...) Fortunately, the Bradley-Terry model \textit{depends only on the difference of rewards} between two completions (...) Substituting the reparameterization (...) into the preference model, \textit{the partition function cancels}, and we can express the human preference probability in terms of only the optimal policy $\pi^*$ and reference policy $\pi_\text{ref}$. \\
\midrule
\textbf{Inference probe} & According to the paper `Direct Preference Optimization: Your Language Model is Secretly a Reward Model', the expensive quantity from the optimal-policy form that becomes unnecessary to estimate once DPO rewrites the Bradley-Terry preference model in terms of policies is \textbf{the partition function}. \\
\bottomrule
\end{tabular}
}}
\caption{\small{The \textit{factual} probe extracts questions from knowledge-bearing sentences, converts them into statements with answers at the end, and adds context. Notice how the italicized words in the original sentence reappear. The \textit{inference} probe composes knowledge from one or multiple sentences (see related sentences) to infer information that was not explicitly stated before. In this case, it integrates three sentences---that the partition function $Z(x)$ is expensive to estimate, that the Bradley-Terry model depends only on reward differences, and that substituting the reparameterization makes the partition function cancel---to infer that the partition function is the quantity that no longer needs to be estimated. The \textbf{target span} of each probe is bolded.}}
\label{tab1}
\vspace{-5mm}
\end{table*}
